\chapter{Project Roadmap}\label{ch:tb-roadmap}

\section*{Scope of this chapter}
\addcontentsline{toc}{section}{Scope of this chapter}

This chapter is the roadmap for the \textbf{Trial Balance review and AI
augmentation} project. It is a project-level artefact, not a programme
plan. Other SAP OSS Finance projects (Simula, Arabic AP, Regulations)
run on the same platform, workspaces, and backends, but each owns its
own roadmap, business case, and go/no-go decision. This chapter adopts
the MoSCoW prioritisation and sequence-based critical path that the
four projects share so that artefacts can be compared at a glance; it
does not assert coupling between their go-live moments.

\paragraph{Sequence, not dates.}
The milestones and dependency graph below express \emph{order} and
\emph{prerequisite} relationships, not calendar dates. This is
deliberate. Implementation teams must know what depends on what, and
where the schedule can absorb slippage (the Should/Could items); a
date column would harden targets into commitments they were never
meant to be. Dates, when set, live in the project steering packet and
the ticketing system, not in this specification.

\section{MoSCoW classification}\label{sec:tb-moscow}

Every milestone below carries one of four priorities. The mapping is
contractual: a milestone's priority determines what happens when it is
dropped, de-scoped, or slips.

\begin{table}[htbp]
\centering
\caption{MoSCoW prioritisation (shared across the four SAP OSS Finance
projects).}
\label{tab:tb-moscow}
\footnotesize
\begin{tabularx}{\linewidth}{@{}l l X@{}}
\toprule
\textbf{Priority} & \textbf{Meaning} & \textbf{Allowance} \\ \midrule
\textbf{Must} (\textbf{M}) & Non-negotiable for v1.0. Without this milestone, the release does not ship. & None. Missing a Must halts the release; see the kill-switch in \cref{sec:tb-kill-switch}. \\
\textbf{Should} (\textbf{S}) & Important but not vital. Ship without it only with sponsor sign-off; document the deviation in the business-case refresh. & May slip by one rollout wave; two consecutive slips escalate to GFAIC. \\
\textbf{Could} (\textbf{C}) & Desirable enhancement. Included if Must and Should items are on track. & May be dropped to v1.1 at project steering's discretion. \\
\textbf{Won't} (\textbf{W}) & Explicitly out of v1.0 scope. Listed so the boundary is legible. & Not eligible for v1.0 planning at all. \\
\bottomrule
\end{tabularx}
\end{table}

\noindent
The critical path of this project is the chain of Must items in
\cref{tab:tb-milestones}; the Should and Could items are where the
schedule and scope can absorb delivery risk without a release slip.

\section{Dependency graph}\label{sec:tb-dependency-graph}

\begin{figure}[htbp]
\centering
\begin{tikzpicture}[
  node distance=5mm and 10mm,
  must/.style={draw,rounded corners=2pt,align=center,font=\scriptsize,
               fill=tbamber!20,draw=tbamber!70,
               minimum height=7mm,minimum width=30mm},
  should/.style={draw,rounded corners=2pt,align=center,font=\scriptsize,
               fill=tbblue!12,draw=tbblue!50,
               minimum height=7mm,minimum width=30mm},
  could/.style={draw,rounded corners=2pt,align=center,font=\scriptsize,
               fill=tbgreen!12,draw=tbgreen!50,dotted,
               minimum height=7mm,minimum width=30mm},
  dep/.style={-{Latex[length=2mm]},thick},
  soft/.style={-{Latex[length=2mm]},thick,dashed}
]
  \node[must] (m01) {M01 TB-ENG-001 \\ \scriptsize dashboard shell};
  \node[must,below=of m01] (m02) {M02 TB-ENG-002 \\ \scriptsize anomaly engine};
  \node[must,below=of m02] (m03) {M03 TB-ENG-003 \\ \scriptsize OAuth scopes};
  \node[must,right=14mm of m02] (m04) {M04 REG-INTG-TB-001 \\ \scriptsize identity wiring};
  \node[must,right=14mm of m04] (m05) {M05 TB-BETA-001 \\ \scriptsize HK beta};
  \node[must,below=of m05] (m06) {M06 TB-PARALLEL-001 \\ \scriptsize parallel-run gate};
  \node[must,below=of m06] (m07) {M07 TB-GA-001 \\ \scriptsize GA sign-off};
  \node[should,right=14mm of m07] (s01) {S01 TB-WAVE-002 \\ \scriptsize wave 2 (3 LEs)};
  \node[should,above=of s01] (s02) {S02 TB-WAVE-003 \\ \scriptsize wave 3 (12 LEs)};
  \node[could,above=of s02] (c01) {C01 TB-V2-KICKOFF \\ \scriptsize v2.0 ML kickoff};

  \draw[dep] (m01) -- (m02);
  \draw[dep] (m01) -- (m03);
  \draw[dep] (m02) -- (m04);
  \draw[dep] (m03) -- (m04);
  \draw[dep] (m04) -- (m05);
  \draw[dep] (m05) -- (m06);
  \draw[dep] (m06) -- (m07);
  \draw[dep] (m07) -- (s01);
  \draw[dep] (s01) -- (s02);
  \draw[soft] (m07) -- (c01);
\end{tikzpicture}
\caption{Trial Balance dependency graph. Amber nodes are Must
(critical path); blue are Should (allowance); dotted green are Could
(discretionary). Solid arrows are hard prerequisites; dashed arrows
indicate conditional starts (the target is unlocked when the upstream
completes \emph{and} the v2.0 trigger fires).}
\label{fig:tb-dependency-graph}
\end{figure}

\section{Project milestones}\label{sec:tb-milestones}

\begin{table}[htbp]
\centering
\caption{Trial Balance v1.0 project milestones in dependency order.
Priority: M=Must (critical path), S=Should, C=Could, W=Won't.}
\label{tab:tb-milestones}
\footnotesize
\begin{tabularx}{\linewidth}{@{}l l l l X@{}}
\toprule
\textbf{ID} & \textbf{Prereq} & \textbf{Pri.} & \textbf{Ticket} & \textbf{Milestone} \\ \midrule
M01 & ---       & M & TB-ENG-001         & Dashboard shell in UAT; all 28 workflow states reachable on the HK fixture. \\
M02 & M01       & M & TB-ENG-002         & Statistical anomaly engine passes the fixtures in \cref{ch:controls-engine}. \\
M03 & M01       & M & TB-ENG-003         & OAuth 2.0 scopes enforced end-to-end (\cref{sec:dashboard-api-auth}). \\
M04 & M02, M03  & M & REG-INTG-TB-001    & Identity-attribution wiring (reserve/complete) live in UAT. \\
M05 & M04       & M & TB-BETA-001        & Beta cycle: HK LE only, AI parallel to manual, no business users yet. \\
M06 & M05       & M & TB-PARALLEL-001    & Beta parallel-run gate passes. \\
M07 & M06       & M & TB-GA-001          & GA sign-off by Head of Africa GFS; HK goes live. \\
M08 & M07       & M & TB-BC-REFRESH      & Business-case refresh and compliance attestation (first cycle after GA; then annual). \\
S01 & M07       & S & TB-WAVE-002        & Rollout wave 2: three further LEs onboarded via entity-params (\cref{sec:params-234}). \\
S02 & S01       & S & TB-WAVE-003        & Rollout wave 3: cumulative 12 LEs live. \\
S03 & S02       & S & TB-WAVE-N          & Subsequent waves through full 234-LE coverage (repeated application of the onboarding runbook). \\
C01 & M07 + v2.0 trigger & C & TB-V2-KICKOFF & v2.0 ML anomaly programme kickoff (Chapter~\ref{ch:future-ml}). \\
C02 & M07       & C & TB-COMMENTARY-TUNE & Quarterly commentary-prompt retraining from accept/edit telemetry. \\
W01 & ---       & W & TB-V1-ML           & ML-based anomaly detection is \emph{not} in v1.0 scope (see \cref{ch:future-ml}). \\
W02 & ---       & W & TB-V1-FULL-234     & Full 234-LE rollout is \emph{not} in v1.0 scope; scope for v1.0 ends at wave 3. \\
\bottomrule
\end{tabularx}
\end{table}

\noindent
The Must chain (M01 $\to$ M02/M03 $\to$ M04 $\to$ M05 $\to$ M06 $\to$
M07 $\to$ M08) is the critical path. The Should items (S01--S03) are
the primary \emph{schedule allowance}: if Must items run hot, rollout
waves can slip without affecting v1.0 sign-off. The Could items are
the \emph{scope allowance}: they may be cut to v1.1 without renegotiation.

\section{Launch tiers and exit gates}\label{sec:tb-launch-tiers}

The project walks the four-tier ladder shared across SAP OSS Finance
projects. A tier is exited only when the named exit gate is satisfied;
slippage requires sponsor sign-off.

\begin{table}[htbp]
\centering
\caption{Trial Balance launch tiers. The exit gate defines the MoSCoW
items that must be closed to exit the tier.}
\label{tab:tb-launch-tiers}
\footnotesize
\begin{tabularx}{\linewidth}{@{}l l X@{}}
\toprule
\textbf{Tier} & \textbf{Scope} & \textbf{Exit gate (references MoSCoW IDs)} \\ \midrule
Alpha & Engineering only, DEV + UAT environments, no business users. & M01, M02, M03 closed; all acceptance criteria in \cref{ch:controls} pass on the HK golden fixture; reproducibility of the variance engine confirmed on two consecutive runs. \\
Beta  & One LE (Hong Kong), parallel to the manual process. & M04, M05, M06 closed; parallel-run comparison meets thresholds for one full cycle; no open SEV1/SEV2; audit coverage $=$ 100\,\% of in-scope actions. \\
GA (limited) & Three LEs (HK plus wave 2). & M07 closed; Beta gate sustained for two consecutive cycles; S01 closed; entity-params onboarding runbook exercised end-to-end; kill-switch green. \\
Full rollout & All LEs in scope for v1.0 (through wave 3). & S02 closed; two quarterly reviews sustained; M08 complete; business-case refresh complete; EUC retirement plan (\cref{sec:euc-retirement-plan}) at ``dashboard authoritative'' for all three EUCs. \\
\bottomrule
\end{tabularx}
\end{table}

\section{Kill-switch conditions}\label{sec:tb-kill-switch}

The kill-switch is a pre-committed rulebook: a tripped condition enters
the action column without further debate. GFAIC and sponsor are
notified within 24 hours. The conditions are evaluated continuously
(on every cycle) and do not carry dates.

\begin{table}[htbp]
\centering
\caption{Trial Balance kill-switch.}
\label{tab:tb-kill-switch}
\footnotesize
\begin{tabularx}{\linewidth}{@{}l l X@{}}
\toprule
\textbf{Condition} & \textbf{Evaluation} & \textbf{Action} \\ \midrule
Parallel-run disagreement rate $>$\,10\,\% on three consecutive cycles & per cycle & \textbf{Abandon v1.0}; fall back to manual; re-scope v1.1. \\
Cycle time from \stateid{EXTRACT\_TB} to \stateid{SEND\_FORTM} $>$\,manual baseline for two cycles & per cycle & \textbf{Redesign}: halt rollout, fix latency, re-enter Beta. \\
Accept-without-edit rate of AI commentary $<$\,40\,\% for two cycles & per cycle & \textbf{Redesign prompt layer}; commentary drafts gated behind a feature flag. \\
Audit-coverage gap (missed \texttt{reserve/complete} pair) detected in production & continuous & \textbf{Pause GA onboarding}; incident-review before next wave. \\
Any MAS MGF finding attributable to TB specifically & continuous & \textbf{Halt rollout}; remediate before next wave. \\
Shared-platform dependency in \cref{tab:tb-shared-deps} declared red by its owner project & continuous & \textbf{Pause downstream Must}; re-plan at next steering. \\
\bottomrule
\end{tabularx}
\end{table}

\section{Shared-platform dependencies}\label{sec:tb-shared-deps}

This project depends on four shared-platform components. Each is owned
by a sibling project (Regulations) or by shared infra, and is tracked
here only as a dependency, not a commitment this project can
unilaterally adjust. Sequence is expressed in terms of the Must item
that first consumes the dependency, not in calendar time.

\begin{table}[htbp]
\centering
\caption{Shared-platform dependencies.}
\label{tab:tb-shared-deps}
\footnotesize
\begin{tabularx}{\linewidth}{@{}l l l X@{}}
\toprule
\textbf{Component} & \textbf{Owner} & \textbf{First consumed by} & \textbf{Impact if late} \\ \midrule
Audit service (reserve/complete, immutable storage) & Regulations project & M04 (REG-INTG-TB-001) & No Beta gate possible without this; TB GA slips in lockstep with the audit service. \\
Unified schema registry (\path{docs/schema/}) & Shared infra & M01 & Already \texttt{migration.status=complete}; no remaining risk. \\
vLLM hosting (\texttt{vllm-only} policy) & Shared platform & M02 & Graceful-degradation ladder defined in the Regulations spec provides fallback; slip tolerable. \\
Review dashboard infra (SAPUI5 + BTP destination) & Shared platform & M01 & Degrade to read-only evidence mirror on Shared Drive; EUC retirement timeline reopens. \\
\bottomrule
\end{tabularx}
\end{table}

\noindent
Shared-platform dependencies are surfaced at monthly steering. If any
upstream owner declares \textbf{amber} or \textbf{red} on a component
in \cref{tab:tb-shared-deps}, the affected downstream Must milestone
is re-planned at the following steering.

\section{Governance and business-case cadence}\label{sec:tb-governance}

\paragraph{Project governance.}
The Trial Balance project reports monthly to a project steering
(Process Owner, Head of Africa GFS, AI Engineer) and quarterly to the
cross-project GFAIC (\cref{ch:controls}). Monthly steering accepts or
rejects amber/red signals on Must items; GFAIC reviews compliance
posture and any cross-project dependencies.

\paragraph{Business-case refresh (M08).}
The 2025-12-14 business case that authorises this project is refreshed
on the following rhythm:

\begin{itemize}
  \item \textbf{Mandatory refresh}: at the first cycle after GA, and
        annually thereafter (ticket TB-BC-REFRESH, priority M in
        \cref{tab:tb-milestones}). Refresh packet includes volume
        assumptions, unit economics,
        FTE leverage realised vs.\ forecast, and a deviation statement
        against the 2025-12-14 baseline.
  \item \textbf{Trigger-based refresh}: any of
        \begin{itemize}
          \item LE count in scope changes by $\ge\,20$\,\% from forecast
                (either direction);
          \item per-cycle run cost changes by $\ge\,30$\,\% (vLLM
                licensing, GPU price, variance-engine throughput);
          \item new regulatory obligation materially affecting scope
                (e.g.\ new MAS MGF pillar, additional jurisdiction);
          \item any kill-switch in \cref{tab:tb-kill-switch} is tripped.
        \end{itemize}
  \item \textbf{Owner}: business sponsor (TB-REQ-001 author,
        S.~Badrinath, or delegate). Approval routed through the
        project steering and filed at
        \path{docs/tb/business-case/refresh-history.yaml}.
\end{itemize}
